\section{Arithmetic estimates for a finite-order Eisenstein-field cutoff}
\label{sec:hecke-arithmetic-attempt}

This section treats the arithmetic part of a radial Hankel cutoff for
$K=\mathbb Q(\sqrt{-3})$.  It distinguishes the selection at concentration
level $1/2$ from the explicitly different selection at level
$1-\Lambda^{-M}$.  The latter admits a crude arithmetic Gram estimate
and a rapidly vanishing normalized theta tail.  These statements do not
by themselves identify the matching local upper limit.
Section~\ref{sec:hecke-original-upper} proves that remaining
local comparison.

Write $\mathcal O=\mathbb Z[\omega]$, with $\omega=e^{2\pi i/3}$,
and let $v=\sqrt3/2$ be its Euclidean covolume.  Fix a primitive
finite-order Hecke character $\chi$ of conductor $\mathfrak f$.  Put
\[
 Q_\chi=3N\mathfrak f,\qquad
 c=\frac{2\pi}{\sqrt{Q_\chi}},\qquad
 r_K=\frac{\pi}{3\sqrt3},\qquad
 \delta_\chi=\mathbf1_{\chi=1}.
\]
All constants may depend on this fixed character.  The raw Hilbert
space is $L^2((0,\infty),du)$, and
\[
 (H_c\phi)(y)=c\int_0^\infty J_0(2c\sqrt{uy})\phi(u)\,du,
 \qquad I\phi=\int_0^\infty\phi(u)\,du,
\]
\[
 E_\chi\phi(x)=\sqrt x\sum_{\mathfrak a}
       \chi(\mathfrak a)\phi(N\mathfrak a\,x)
       -\delta_\chi r_K x^{-1/2}I\phi.
\]
The multiplicative variable $x$ becomes logarithmic time $t=\log x$.
The raw dilation Sobolev space $H_D^s$ is ordinary $H^s(\mathbb R)$
under $U\phi(t)=e^{t/2}\phi(e^t)$.

\subsection{The lattice theta identity and a Riesz estimate}

\begin{proposition}[Finite-order lattice normalization]
\label{prop:hecke-lattice-poisson}
There is a constant $\varepsilon_\chi$ of modulus one such that
\[
 E_\chi\phi(x)=\varepsilon_\chi
                  E_{\bar\chi}(H_c\phi)(x^{-1}).
\]
For the principal character $\varepsilon_1=1$ and
\[
 H_c\phi(0)=cI\phi,\qquad I(H_c\phi)=\phi(0)/c.
\]
Define the norm-counting discrepancy
\[
 A_\chi(X)=\sum_{N\mathfrak a\leq X}\chi(\mathfrak a)
                          -\delta_\chi r_KX.
\]
Its first integral has the absolutely convergent expansion
\begin{equation}\label{eq:hecke-riesz}
 R_\chi(X):=\int_0^X A_\chi(y)\,dy
 =-\frac{\delta_\chi X}{6}
 +\frac{\varepsilon_\chi X}{c}
   \sum_{\mathfrak a}\frac{\bar\chi(\mathfrak a)}{N\mathfrak a}
                  J_2(2c\sqrt{XN\mathfrak a}).
\end{equation}
Consequently
\[
 R_\chi(X)=-\delta_\chi X/6+O_\chi(X^{3/4})\quad(X\geq1).
\]
Also $A_\chi(X)=O_\chi(X^{1/3})$ for $X\geq1$.
\end{proposition}

\begin{proof}
The class number is one, and each nonzero principal ideal has six
generators.  Write $\mathfrak f=(F)$.  The ideal character is represented
by a multiplicative residue character on $\mathcal O/(F)$ which is
trivial on the six units and zero on nonunits.  Hence the ideal theta
sum is one sixth of the corresponding weighted lattice sum.  The dual
lattice of $\mathcal O$ is $(i/v)\mathcal O$; the dual of
$F\mathcal O$ is $(i/(v\bar F))\mathcal O$.
Finite Fourier transformation of a primitive residue character gives
its conjugate times a Gauss sum of modulus $\sqrt{N\mathfrak f}$.
For a radial function $\phi(|z|^2x)$, Euclidean Fourier transformation
gives
\[
 \widehat{\phi(|\cdot|^2x)}(\xi)
   =\frac\pi x\int_0^\infty
             J_0(2\pi|\xi|\sqrt{u/x})\phi(u)\,du.
\]
Since $|\xi|=c\sqrt{N\mathfrak a}/\pi$ on the relevant dual lattice,
the Gauss factor and covolume combine to give
$\varepsilon_\chi x^{-1}H_c\phi(N\mathfrak a/x)$.
This proves the nonprincipal identity.  In the principal case the
two lattice origins give
\[
 6\sum_{\mathfrak a}\phi(N\mathfrak a x)+\phi(0)
 =x^{-1}\left(6\sum_{\mathfrak a}
 H_c\phi(N\mathfrak a/x)+cI\phi\right).
\]
Here $c/6=r_K$.  Subtracting the indicated mean proves the corrected
identity.  The two identities at zero follow from $H_c^2=I$.

Apply the same lattice formula to $(X-|z|^2)_+$.  At a nonzero
frequency $\xi$ its Fourier transform is
\[
 \frac{X}{\pi|\xi|^2}J_2(2\pi\sqrt X|\xi|),
\]
as follows by integrating $2\pi r(X-r^2)J_0(2\pi r|\xi|)$.
At zero the integral is $\pi X^2/2$.  The latter is the main term
$r_KX^2/2$ after division by the six units and the covolume.
Removing the original lattice origin contributes $-X/6$ in the
principal case.  Grouping the remaining dual points into ideals
proves \eqref{eq:hecke-riesz}.  Since
$J_2(z)=O(z^{-1/2})$ for $z\geq1$, each summand is bounded by a
constant times $X^{3/4}(N\mathfrak a)^{-5/4}$.  The latter series
converges, because the number of ideals of norm at most $Y$ is
$O_K(Y)$.

For completeness the weaker pointwise discrepancy estimate follows
from elementary smoothing.  Count points of each of the finitely
many residue-class lattices in a disk of radius $T=\sqrt X$.
Convolve indicators of the disks of radii $T\pm\eta$ with a
nonnegative smooth mollifier supported in a disk of radius $\eta$.
They respectively majorize and minorize the sharp disk indicator.
Their zero Fourier coefficients differ from the disk area by
$O(T\eta+\eta^2)$.  At nonzero dual frequencies the disk transform
is $O(T^{1/2}|\xi|^{-3/2})$.  The mollifier contributes
$O_N((1+\eta|\xi|)^{-N})$, so summation over a two-dimensional
dual lattice gives $O(T^{1/2}\eta^{-1/2})$.
Choose $\eta=T^{-1/3}$.  This proves discrepancy $O(T^{2/3})$
for each translated lattice.  Taking the finite character-weighted
combination proves $A_\chi(X)=O_\chi(X^{1/3})$.
\end{proof}

\begin{proposition}[Improper Mellin integrals at sharp endpoints]
\label{prop:hecke-improper-mellin}
Let $\phi$ be smooth on $[0,A]$ and zero outside that interval; a
jump at $A$ is allowed.  For $0<\Re s<1$,
\begin{equation}\label{eq:hecke-improper-mellin}
 \lim_{\eta\downarrow0}\int_\eta^\infty
 E_\chi\phi(x)x^{s-1/2}\frac{dx}{x}
 =L_K(s,\chi)\int_0^A\phi(u)u^{s-1}\,du.
\end{equation}
The convergence is locally uniform in this strip, and any fixed
number of derivatives in $s$ may be taken.  Absolute convergence at
zero is not asserted.
\end{proposition}

\begin{proof}
First take $\phi=1_{(0,a)}$.  Its corrected theta image is
$\sqrt x A_\chi(a/x)$.  For a nonprincipal character, integration
by parts against $R_\chi$ gives, uniformly on compact subsets of
$\Re s>-1/4$,
\[
 \int_Y^\infty A_\chi(y)y^{-s-1}\,dy
 =-R_\chi(Y)Y^{-s-1}
 +(s+1)\int_Y^\infty R_\chi(y)y^{-s-2}\,dy
 =O(Y^{-\Re s-1/4}).
\]
For the principal character separate the term $-y/6$ in $R_1$;
its contribution to the displayed integral is $O(Y^{-\Re s})$.
Thus the integral converges for every $\Re s>0$.  The same estimates
with fixed powers of $\log y$ prove locally uniform convergence of
its derivatives.  The change of variables $y=a/x$ is legitimate in
each truncated integral.  For a nonprincipal character,
Stieltjes summation first gives
$L_K(s,\chi)/s=\int_1^\infty A_\chi(y)y^{-s-1}\,dy$
when $\Re s>1$.  The preceding locally uniform convergence
continues this identity to $\Re s>0$.  Since $A_\chi(y)=0$
for $0<y<1$, multiplication by $a^s$ proves the disk identity.

For the principal character there is no common region of absolute
convergence for the whole corrected disk integral and the original
Dirichlet series.  Instead, Stieltjes summation for $\Re s>1$ gives
\[
 \frac{L_K(s,1)}s
 =\int_1^\infty A_1(y)y^{-s-1}\,dy+\frac{r_K}{s-1}.
\]
The integral on the right continues holomorphically to $\Re s>0$
by the Riesz estimate already proved.  On $0<\Re s<1$ one has
$A_1(y)=-r_Ky$ for $0<y<1$, and therefore
\[
 \frac{r_K}{s-1}=\int_0^1 A_1(y)y^{-s-1}\,dy.
\]
Combining these identities proves
$L_K(s,1)/s=\int_0^\infty A_1(y)y^{-s-1}\,dy$
in this strip, and hence the principal disk identity after multiplying
by $a^s$.  The restriction $\Re s<1$ is exactly the convergence
condition for the mean-subtraction tail at $x=\infty$.

For a general sharp profile use the exact representation
\[
 \phi(u)=\phi(A)1_{u<A}
                -\int_0^A\phi'(a)1_{u<a}\,da.
\]
The omitted lower Mellin integral for a nonprincipal disk is
$O(\eta^{\Re s+1/4}a^{-1/4})$; the same estimate remains valid
when $a<\eta$ by using the complete disk integral.  The principal
case has in addition $O(\eta^{\Re s})$.  Fixed derivatives only
insert fixed powers of $|\log\eta|$ and $|\log a|$.
Because $a^{-1/4}(1+|\log a|)^j$ is integrable at zero, these bounds
justify integration against $\phi'(a)$, the interchange of the
truncated integrals, and passage to their limits.  They also justify
all fixed Mellin jets.  This proves the proposition.
\end{proof}

The improper Mellin statement must not be replaced by a convergence
claim for the ordinary Dirichlet-series tail.  In fact the exact
analytically regularized tail is
\begin{equation}\label{eq:hecke-tail-endpoint}
 L_K(s,\chi)-\sum_{N\mathfrak a\leq X}
    \frac{\chi(\mathfrak a)}{(N\mathfrak a)^s}
 +\delta_\chi\frac{r_KX^{1-s}}{1-s}
 =-A_\chi(X)X^{-s}
           +s\int_X^\infty A_\chi(y)y^{-s-1}\,dy.
\end{equation}
The endpoint discrepancy survives, even though the integral on the
right has the stronger Riesz bound.

\subsection{A graph domain and the crude Gram estimate}

We use the explicit analytic input
\begin{equation}\label{eq:hecke-subconvex-input}
 |L_K(1/2+i\tau,\chi)|\leq C_\chi(2+|\tau|)^{3/8}.
\end{equation}
For a cuspidal finite-order theta series, including the cubic
character of conductor $(269)$, this follows from the fixed-form
Weyl bound of \citet[Theorem~1.1, p.~2]{booker2019}, which gives
$|\tau|^{1/3}\log|\tau|$ for a primitive holomorphic cusp form of
arbitrary level.  The theorem states no exclusion of weight one.
In a reducible theta sector the $L$-function is a product of two
fixed Dirichlet $L$-functions.  The all-conductor Weyl bound of
\citet[Theorem~1.1, p.~1]{petrowyoung2023} gives exponent
$1/6+\epsilon$ for each factor, hence $1/3+2\epsilon<3/8$
after choosing a fixed sufficiently small $\epsilon$.
Thus it gives the same conclusion.  A proof using this section must retain
this subconvexity input; degree-two convexity alone reaches the
unsuitable endpoint exponent $1/2$.

Plancherel and \eqref{eq:hecke-subconvex-input} give
\begin{equation}\label{eq:hecke-graph}
 \|E_\chi\phi\|_2\leq C_\chi\|\phi\|_{H_D^{3/8}}.
\end{equation}
This defines the closed theta operator on that Sobolev space.
For sharp time-supported smooth profiles it agrees almost
everywhere with the displayed finite theta sum: smooth the endpoint
in $H_D^{3/8}$, use \eqref{eq:hecke-graph} for $L^2$ convergence,
and observe convergence of the finite sum at every $x$ outside the
countable set of endpoint ratios.  The principal means converge by
ordinary $L^2$ approximation on a fixed compact interval.
The preceding improper Mellin proposition gives the full complex
Mellin identity and its jets for these same profiles.

Set
\[
 P_\Lambda=1_{(0,\Lambda)},\qquad
 B_\Lambda=H_cP_\Lambda H_c,\qquad
 A_\Lambda=P_\Lambda B_\Lambda P_\Lambda.
\]
Let $\mathcal R_\Lambda$ be any spectral subspace of $A_\Lambda$
with eigenvalues at least $\nu>0$, with the two constraints
$I\phi=\phi(0)=0$ imposed in the principal sector.  We use the
archimedean mass estimate
\begin{equation}\label{eq:hecke-core-mass}
 \int_{d/(c^2\Lambda)}^\Lambda|\phi(u)|^2du
                         \geq\tfrac12\|\phi\|_2^2
 \qquad(\phi\in\mathcal R_\Lambda).
\end{equation}
It follows directly by extending $\phi$ to a Hankel-band vector
of norm at most $\nu^{-1/2}\|\phi\|$ and using
$|w(u)|^2\leq c^2\Lambda\|w\|^2$ near zero.

\begin{proposition}[Crude Hecke Gram control]
\label{prop:hecke-crude-gram}
For a raw orthonormal basis of $\mathcal R_\Lambda$, let $G_\Lambda$
be the Gram matrix of its theta images.  Then
\[
 \|G_\Lambda^{-1}\|\leq C_{\chi,\nu}\Lambda^2.
\]
No cancellation estimate for the ideal M\"obius function is required.
\end{proposition}

\begin{proof}
Put $a=d/(c^2\Lambda)$ and $N=\lceil\Lambda/a\rceil$.
For $f=U\phi$ and $e=E_\chi\phi$, finite ideal M\"obius inversion
gives exactly, for $\log a\leq t\leq\log\Lambda$,
\[
 f(t)=\sum_{N\mathfrak a\leq N}
  \frac{\mu_K(\mathfrak a)\chi(\mathfrak a)}{\sqrt{N\mathfrak a}}
                         e(t+\log N\mathfrak a).
\]
Expansion groups terms by the product ideal $\mathfrak b$ and uses
$\sum_{\mathfrak a\mid\mathfrak b}\mu_K(\mathfrak a)
=\mathbf1_{\mathfrak b=\mathcal O}$.  The support of $\phi$ removes
all products beyond $\Lambda e^{-t}\leq N$.  In the principal sector
$I\phi=0$ removes the correction term, so the same finite identity
holds.  Its translation operator has $L^2$ norm at most
\[
 \sum_{N\mathfrak a\leq N}\frac1{\sqrt{N\mathfrak a}}
                         \leq C_K\sqrt N\leq C_\chi\Lambda,
\]
by the elementary ideal count $O_K(Y)$ and partial summation.
The core mass estimate therefore yields
$\|\phi\|_2\leq C_\chi\Lambda\|E_\chi\phi\|_2$.
Squaring is precisely the asserted Gram inequality.
\end{proof}

\subsection{Deep concentration and unconditional lower estimates}

For a fixed $M>0$ make the explicit choice
\[
 \Pi_{\Lambda,M}
       =\mathbf1_{[1-\Lambda^{-M},1]}(A_\Lambda).
\]
The principal raw column space is its intersection with
$\ker I\cap\ker\operatorname{ev}_0$.  This is a different cutoff
from $\mathbf1_{[1/2,1]}(A_\Lambda)$.  Let $E_\Lambda$ denote the
theta column operator, $Z_\Lambda=E_\Lambda G_\Lambda^{-1/2}$,
$Q_\Lambda=Z_\Lambda Z_\Lambda^*$, and
$S_L=1_{[-L,L]}$, $L=\log\Lambda$.

\begin{theorem}[Deep-cut theta leakage]
\label{thm:hecke-deep-leakage}
Assume the graph bound \eqref{eq:hecke-graph} and the archimedean
bound $\|\psi\|_{H_D^s}\leq C_s\Lambda^s\|\psi\|_2$ on the
selected raw space for $s<1/2$.  More generally let the graph exponent
be $\sigma<s<1/2$.  Then
\begin{equation}\label{eq:hecke-deep-leakage}
 \|(1-S_L)Z_\Lambda\|_2
 \leq C\Lambda^{3/2+\sigma-(M/2)(1-\sigma/s)}.
\end{equation}
In particular, $\sigma=3/8$, $s=7/16$, and $M=56$ give
$O(\Lambda^{-17/8})$.  For this choice,
\[
 \operatorname{Tr}((Q_\Lambda-S_LQ_\Lambda S_L)V(h))\longrightarrow0
 \qquad(h\in C_c^\infty(\mathbb R)).
\]
\end{theorem}

\begin{proof}
Write $F_\Lambda=P_\Lambda H_cP_\Lambda$, so $A_\Lambda=F_\Lambda^2$.
On the unconstrained deep selected span put
$W=B_\Lambda A_\Lambda^{-1}$, so $P_\Lambda W=I$.  Spectral calculus
gives
\[
 \|(1-P_\Lambda)W\psi\|_2^2
 =\langle\psi,(A_\Lambda^{-1}-I)\psi\rangle
 \leq C\Lambda^{-M}\|\psi\|_2^2.
\]
Furthermore $H_cW\psi=F_\Lambda^{-1}\psi$ belongs to the same
selected time-supported span.  Hankel transformation reverses the
dilation generator and preserves its Sobolev norms.  Thus both
$W\psi$ and $\psi$ have $H_D^s$ norm $O(\Lambda^s)\|\psi\|$.
Interpolation gives
\[
 \|(1-P_\Lambda)W\psi\|_{H_D^\sigma}
 \leq C\Lambda^{\sigma-(M/2)(1-\sigma/s)}\|\psi\|_2.
\]
The selected dimension is $O(\Lambda)$, since it is bounded by
twice $\operatorname{Tr}A_\Lambda$ for large $\Lambda$.  Sum the
squared estimate over a raw orthonormal basis and apply the graph
bound.  The identity $E_\chi H_c=\varepsilon_\chi J E_{\bar\chi}$,
where $(Jf)(t)=f(-t)$, extends to $H_D^\sigma$ because $H_c$
preserves that space and both theta maps satisfy the graph bound.
Apply it to $H_cW\psi=F_\Lambda^{-1}\psi$, which is time supported.
In the principal case its mean is
$I(H_cW\psi)=W\psi(0)/c=\psi(0)/c=0$.
It follows directly that the lower theta tail of $W\psi$ is zero.
The upper theta tail of $\psi$ is zero,
using $I\psi=0$ in the principal case.  Consequently
$\|(1-S_L)E_\Lambda\|_2$ is bounded by the theta norm just estimated.
Finally $\|G_\Lambda^{-1/2}\|=O(\Lambda)$ proves
\eqref{eq:hecke-deep-leakage}.

For the trace statement use
$Q-S_LQS_L=(1-S_L)Q+S_LQ(1-S_L)$ and Hilbert--Schmidt
Cauchy--Schwarz.  Since $\|Z_\Lambda\|_2=O(\Lambda^{1/2})$,
the absolute trace difference is
$O(\|V(h)\|\Lambda^{1/2}\|(1-S_L)Z_\Lambda\|_2)=o(1)$.
\end{proof}

\begin{proposition}[All zero jets for the deep cutoff]
\label{prop:hecke-deep-zero-jets}
Let $M>4/3$ and let $\mathcal V_L$ be the span of the truncated
exponential-polynomial jets associated with any fixed finite packet
of nontrivial zeros of $L_K(s,\chi)$, including their multiplicities.
Then, for some fixed exponent $C$ depending on the packet,
\[
 \|Q_\Lambda P_{\mathcal V_L}\|_2^2
                       =O(\Lambda^{4/3-M}(1+L)^C)=o(1).
\]
The resulting finite-packet lower bound is the harmonic-measure
bound: a zero $\rho=\beta+i\gamma$ contributes multiplicity times
the point mass at $\gamma$ when $\beta=1/2$, and otherwise the
Poisson measure of width $|\beta-1/2|$.  No zero-location theorem
is used.
\end{proposition}

\begin{proof}
For a zero $s_0=\beta+i\gamma$ and $j<m_{s_0}$, define the lower
Mellin tail with weight $(\log x+L)^j$.  Proposition
\ref{prop:hecke-improper-mellin} justifies this improper integral.
Its raw kernel is
\[
 H_{s_0,j,\Lambda}(u)=
 \left.(\partial_s+L)^j\left[
 u^{s-1}\left(L_K(s,\chi)
       -\sum_{N\mathfrak a\leq\Lambda u}
            \frac{\chi(\mathfrak a)}{(N\mathfrak a)^s}\right)
 +\delta_\chi\frac{r_K\Lambda^{1-s}}{1-s}
 \right]\right|_{s=s_0}.
\]
The formula follows on a truncated compact raw core by Stieltjes
summation; the Riesz bounds in the preceding proof justify its
limit and all fixed derivatives.  For $u<\Lambda^{-1}$, the zero
and its jets remove the $L$-term.  Thus the kernel is zero there
in a nonprincipal sector, and bounded there in the principal
sector.  For $u\geq\Lambda^{-1}$, apply
\eqref{eq:hecke-tail-endpoint}, using the pointwise bound
$A_\chi(X)=O(X^{1/3})$ for its endpoint term and the Riesz estimate
for its conditionally convergent integral.  The latter use is
necessary when $\beta\leq1/3$.  The derivatives involve scale-free
logarithms, so for every fixed jet
\[
 |H_{s_0,j,\Lambda}(u)|
       \leq C\Lambda^{1/3-\beta}u^{-2/3},
 \qquad
 \|1_{u>\Lambda}H_{s_0,j,\Lambda}\|_2^2
       \leq C\Lambda^{1/3-2\beta}.
\]
The global kernel belongs to raw $L^2$.  First apply its identity
to inverse Hankel transforms of smooth compact band profiles.
These are radial Schwartz functions; the disk-superposition proof
extends to them because their derivatives have rapid decay, so all
the same Riesz majorants remain integrable over $0<a<\infty$.
Poisson makes their lower theta tail zero, with the value constraint
in the principal case.  Such vectors are dense in the fixed band
space, and the principal value constraint can be retained because
evaluation there is continuous.  Continuity of the $L^2$ kernel
therefore proves annihilation on $W\psi$, without an assertion of
absolute Mellin convergence.
Therefore, for a raw orthonormal selected basis,
\[
 \sum_i|\ell_{s_0,j,\Lambda}(\psi_i)|^2
 \leq\|(1-P_\Lambda)W\|^2
         \|1_{u>\Lambda}H_{s_0,j,\Lambda}\|_2^2
 \leq C\Lambda^{-M+1/3-2\beta}.
\]
The inverse Gram bound changes this to
$O(\Lambda^{7/3-M-2\beta})$ for normalized theta columns.
The full Mellin jet vanishes and the upper theta tail is zero.
For $\alpha=\beta-1/2<0$, use endpoint vectors
$1_{[-L,L]}e^{(\alpha-i\gamma)(t+L)}(t+L)^j$; their extra squared
factor $\Lambda^{2\alpha}$ leaves $\Lambda^{4/3-M}$.
For $\alpha>0$, use the right-endpoint normalization and expand
$(t-L)^j$ in powers of $(t+L)$; the power is
$\Lambda^{10/3-M-4\beta}\leq\Lambda^{4/3-M}$, with fixed
logarithmic factors.  Central-line polynomial jets have their
usual polynomial $L$ normalization and satisfy the same bound.
The finite endpoint Gram matrices have nondegenerate limits, so
the estimate passes to the orthogonal projection onto their span.
The finite exponential-packet trace calculation then gives the
stated harmonic measures.  Increasing finite packets gives the
corresponding all-zero lower bound whenever the ordinary zero-count
bound is used to sum the resulting positive measures.
\end{proof}

